% =====================================================================
%  config/preamble.tex
%  ---------------------------------------------------------------
%  Paquetes, geometría, tipografía, encabezados y estilo de tablas.
%  Todo lo aquí usado forma parte de una instalación estándar de
%  TeX Live / MacTeX. No se requieren paquetes exóticos ni servicios web.
% =====================================================================

% ------------------------- Codificación e idioma ---------------------
\usepackage[T1]{fontenc}
\usepackage[utf8]{inputenc}
\usepackage{lmodern}
\usepackage[spanish,mexico,es-nolists,es-noindentfirst]{babel}
% Nota: la opción «mexico» usa punto decimal.
%
% Rótulos: la plantilla Word original habla de «cuadros y figuras», y así
% se citan en el texto. Las versiones recientes de babel-spanish rotulan
% las tablas como «Tabla», por lo que se fija explícitamente «Cuadro»
% para que el rótulo automático y las menciones del texto coincidan.
% Si prefiere «Tabla», cambie las dos líneas siguientes y, con ellas,
% las menciones «Cuadro~\ref{...}» del texto.
\addto\captionsspanish{%
  \renewcommand{\tablename}{Cuadro}%
  \renewcommand{\listtablename}{Índice de cuadros}%
  \renewcommand{\listfigurename}{Índice de figuras}%
  \renewcommand{\contentsname}{Tabla de contenido}%
}

% ---- Símbolos Unicode de uso frecuente ------------------------------
% Permiten que el equipo escriba directamente ×, ÷, ≥, ≤, ≠, → … en el
% texto y en las tablas sin recurrir a modo matemático.
\usepackage{amsmath}
\DeclareUnicodeCharacter{00D7}{\ensuremath{\times}}
\DeclareUnicodeCharacter{00F7}{\ensuremath{\div}}
\DeclareUnicodeCharacter{2264}{\ensuremath{\leq}}
\DeclareUnicodeCharacter{2265}{\ensuremath{\geq}}
\DeclareUnicodeCharacter{2260}{\ensuremath{\neq}}
\DeclareUnicodeCharacter{2192}{\ensuremath{\rightarrow}}
\DeclareUnicodeCharacter{2190}{\ensuremath{\leftarrow}}
\DeclareUnicodeCharacter{2248}{\ensuremath{\approx}}
\DeclareUnicodeCharacter{2212}{\ensuremath{-}}
\DeclareUnicodeCharacter{2022}{\textbullet}

% --------------------------- Tipografía ------------------------------
\usepackage{helvet}
\renewcommand{\familydefault}{\sfdefault}   % aproxima la Arial del Word
\usepackage{setspace}
\onehalfspacing
\usepackage{microtype}

% ---------------------------- Geometría ------------------------------
\usepackage{geometry}
\geometry{a4paper,
  left=2.5cm, right=2.5cm, top=2.4cm, bottom=2.2cm,
  headheight=46pt, headsep=12pt, footskip=34pt}

% ------------------------------ Color --------------------------------
\usepackage{xcolor}
\definecolor{uvazul}{HTML}{1F3864}   % azul institucional para títulos
\definecolor{uvgris}{HTML}{D9D9D9}   % encabezados de tabla
\definecolor{uvgrisc}{HTML}{F2F2F2}  % relleno suave
\definecolor{instrverde}{HTML}{2E6B34}
\definecolor{instrfondo}{HTML}{EDF6EE}
\definecolor{polinaranja}{HTML}{9C5A00}
\definecolor{polfondo}{HTML}{FDF3E3}
\definecolor{ejemazul}{HTML}{1F4E79}
\definecolor{ejemfondo}{HTML}{EAF1F8}
\definecolor{normgris}{HTML}{4A4A4A}
\definecolor{normfondo}{HTML}{F4F4F4}
\definecolor{obsmorado}{HTML}{6B3FA0}  % observación de la revisión (variante anotada)
\definecolor{obsfondo}{HTML}{F3EEFA}
\definecolor{campoamar}{HTML}{FFF2CC}
\definecolor{campotexto}{HTML}{9A3412}  % marcador de campo por rellenar

% ------------------------- Tablas y figuras --------------------------
\usepackage{array}
\usepackage{tabularx}
\usepackage{longtable}
\usepackage{xltabular}   % longtable + columnas X: fichas que pueden partirse entre páginas
\usepackage{multirow}
\usepackage{colortbl}
\usepackage{booktabs}
\usepackage{caption}
\usepackage{graphicx}
\usepackage{float}

\captionsetup{font=small,labelfont=bf,justification=centering,skip=6pt}
\captionsetup[table]{position=above}

% Tipos de columna reutilizables ---------------------------------------
% El \hspace{0pt} inicial permite que TeX divida con guion la PRIMERA
% palabra de la celda (sin él, TeX nunca parte la primera palabra de un
% párrafo, y las columnas estrechas se desbordan).
\newcolumntype{Y}{>{\RaggedRight\arraybackslash\hspace{0pt}}X}                       % flexible
\newcolumntype{C}[1]{>{\Centering\arraybackslash\hspace{0pt}}p{#1}}                  % centrada fija
\newcolumntype{L}[1]{>{\RaggedRight\arraybackslash\hspace{0pt}}p{#1}}                % izquierda fija
\newcolumntype{E}[1]{>{\columncolor{uvgrisc}\bfseries\RaggedRight\arraybackslash\hspace{0pt}}p{#1}} % etiqueta

\renewcommand{\arraystretch}{1.25}
\setlength{\tabcolsep}{5pt}

% ------------------------- Índice (TOC) ------------------------------
\usepackage{tocloft}
\setlength{\cftbeforesecskip}{2pt}
% La numeración de los anexos («ANEXO A») es mucho más ancha que la de
% una sección numérica, por lo que \iniciaranexos ensancha la caja del
% número a partir de ese punto del índice (véase config/commands.tex).

% --------------------------- Listas ----------------------------------
\usepackage{ragged2e}   % \RaggedRight conserva la partición de palabras
\usepackage{enumitem}
\setlist{nosep,topsep=3pt,itemsep=2pt,parsep=0pt,leftmargin=1.4em}

% ------------------------ Encabezado y pie ---------------------------
\usepackage{fancyhdr}
\usepackage{lastpage}

% Los preliminares se numeran en romanos y el cuerpo reinicia en arábigos,
% de modo que «de N» sólo tiene sentido dentro del cuerpo: \pageref{LastPage}
% devuelve el último número arábigo, no el total físico de páginas.
% main.tex activa este interruptor al llamar a \pagenumbering{arabic}.
\newif\ifcontarpaginas

\newcommand{\bandaEncabezado}{%
  \setlength{\tabcolsep}{4pt}%
  \renewcommand{\arraystretch}{1.05}%
  \footnotesize
  \begin{tabularx}{\textwidth}{|>{\raggedright\arraybackslash}X|C{3.2cm}|>{\raggedleft\arraybackslash}X|}
    \hline
    \uvPrograma & \textbf{\uvActividad} & \uvEE \\ \hline
    \multicolumn{3}{|c|}{\uvBandaHeader} \\ \hline
  \end{tabularx}%
}

\newcommand{\bandaPie}{%
  \setlength{\tabcolsep}{4pt}%
  \renewcommand{\arraystretch}{1.05}%
  \footnotesize
  \begin{tabularx}{\textwidth}{|>{\raggedright\arraybackslash}X|>{\raggedleft\arraybackslash}p{5.4cm}|}
    \hline
    \textbf{\uvUniversidad}\ \textperiodcentered\ \uvFacultad
      & Página \thepage\ifcontarpaginas\ de \pageref{LastPage}\fi\quad \docVersionPlantilla \\ \hline
  \end{tabularx}%
}

\fancypagestyle{uvestilo}{%
  \fancyhf{}%
  \renewcommand{\headrulewidth}{0pt}%
  \renewcommand{\footrulewidth}{0pt}%
  \fancyhead[C]{\bandaEncabezado}%
  \fancyfoot[C]{\bandaPie}%
}
\pagestyle{uvestilo}
\fancypagestyle{plain}{\pagestyle{uvestilo}}

% ------------------------ Títulos de sección -------------------------
\usepackage{titlesec}
\titleformat{\section}[block]
  {\normalfont\large\bfseries\color{uvazul}}{\thesection.}{0.6em}{}
\titleformat{\subsection}[block]
  {\normalfont\normalsize\bfseries\color{uvazul}}{\thesubsection}{0.6em}{}
\titleformat{\subsubsection}[block]
  {\normalfont\normalsize\bfseries}{\thesubsubsection}{0.6em}{}
\titlespacing*{\section}{0pt}{14pt}{6pt}
\titlespacing*{\subsection}{0pt}{12pt}{4pt}
\titlespacing*{\subsubsection}{0pt}{10pt}{3pt}

% ---------------------- Cajas de la plantilla ------------------------
\usepackage{etoolbox}
\usepackage[most]{tcolorbox}
\usepackage{comment}

% ----------------------------- Enlaces -------------------------------
\usepackage[hidelinks,bookmarksnumbered,bookmarksopen]{hyperref}
\usepackage{bookmark}

% ------------------------------ TikZ ---------------------------------
\usepackage{tikz}
\usetikzlibrary{positioning,arrows.meta,shapes.geometric,fit,backgrounds,calc,decorations.pathreplacing}

% ------------------- Ajustes finales de composición ------------------
\setlength{\parskip}{6pt}
\setlength{\parindent}{0pt}
\setlength{\emergencystretch}{3em}   % reduce overfull hboxes en texto justificado
\widowpenalty=10000
\clubpenalty=10000
